\begin{question}
Adapt the |TerminatingPartialQueue| class from
Figure~\ref{fig:terminating-queue} so as to use a monitor. 

%%  (If you like, you
%% can change the signature, if it's more convenient for the following part.)

Now adapt the concurrent graph search program from
Figure~\ref{fig:graph-search-conc} so as to use shared variables and a
monitor.
%
Test your code by adapting the word path program from the course website.
\end{question}

%%%%%%%%%%%%%%%%%%%%%%%%%%%%%%%%%%%%%%%%%%%%%%%%%%%%%%%

\begin{answerS}
My code for the terminating partial queue is below.  This maintains, in
|waiting|, a count of the number of threads that are currently waiting to
perform a |dequeue|.  If a thread attempts a |dequeue|, and all the other
threads are waiting, it shuts down the queue, and all the threads return
|None|. 
%
\begin{scala}
class MonitorTerminatingPartialQueue[A](numWorkers: Int){
  /** The queue itself. */
  private val queue = new Queue[A]

  /** The number of threads currently waiting to perform a £dequeue£. */
  private var waiting = 0

  /** Has the queue been shut down? */
  private var done = false

  /** Enqueue £x£. */
  def enqueue(x: A) = synchronized{ 
    if(!done){
      queue.enqueue(x)
      if(waiting > 0) notify()
    }
  }

  /** Attempt to dequeue a value.  
    * @return £None£ if the queue has been shut down, or it the queue is empty
    * and all threads are trying to dequeue. */
  def dequeue(): Option[A] = synchronized{
    if(!done && queue.isEmpty){
      if(waiting == numWorkers-1) shutdown() // System should terminate.
      else{ // This thread needs to wait.
        waiting += 1
        while(queue.isEmpty && !done) wait()
        waiting -= 1
      }
    }
    if(done) None 
    else Some(queue.dequeue())
  }

  /** Shut down this queue. */
  def shutdown() = synchronized{ done = true; notifyAll() }
}
\end{scala}

%%%%%

My code is below.  This is mostly a straightforward adaptation of the earlier
version.  

The previously seen nodes are stored in a sharded set.  Recall that we require
the number of shards to be a power of~$2$; we choose the smallest power of~$2$
that is at least $4 \times \sm{numShards}$.

The variable |result| ends up holding the result of the search.  To avoid
races, we protect this by a |synchronized| block. 

If a worker thread finds a path, and it is the first worker to do so, it sets
|result|, and shuts down the terminating queue so that the other threads
terminate.
%
\begin{scala}
class MonitorGraphSearch[N](g: Graph[N], numWorkers: Int) 
    extends GraphSearch[N](g){
  /** Try to find a path in £g£ from £start£ to a node that satisfies £isTarget£. */
  def apply(start: N, isTarget: N => Boolean): Option[List[N]] = {
    // The number of shards to use: a power of £2£, at least £$4 \times \sm{numShards}$£.
    var numShards = 4; while(numShards < 4*numWorkers) numShards *= 2
    // All nodes seen so far.
    val seen = new MonitorShardedSet[N](numShards); seen.add(start)
    // Queue storing edges and the path leading to that edge.
    val queue = new MonitorTerminatingPartialQueue[(N, List[N])](numWorkers)
    queue.enqueue((start, List(start)))    
    // Variable that ends up holding the result; protected by £synchronized£
    // block.
    var result: Option[List[N]] = None

    // A single worker.
    def worker = thread("worker"){
      var done = false
      while(!done) queue.dequeue() match{
        case Some((n, path)) =>
          for(n1 <- g.succs(n)){
            if(isTarget(n1)) synchronized{
              if(result.isEmpty){ // This is the first thread to find a result.
                result = Some(path :+ n1) // Store result.
                queue.shutdown() // Cause others to terminate.
                done = true
              }
            }
            else if(seen.add(n1)) queue.enqueue((n1, path :+ n1))
          } 
        case None => done = true // Termination case.
      } // end of £while£ loop.
    }

    // Release the workers. 
    run(|| (for(_ <- 0 until numWorkers) yield worker))
    result
  }
}
\end{scala}
\end{answerS}
